\section{Six Evidence-Bounded Research Gaps}

The gaps below are gaps in discrimination and comparability within the audited collection. They are not claims that an experiment has never been performed. Each begins with an abstract-level pattern, identifies the alternative explanations left coupled, and states an observation that could falsify the working explanation. This is Contribution C's problem map applied through Contribution B's minimum record.

\subsection{Gap 1: precursor homogeneity versus powder route}

Three stored abstracts provide complementary anchors: precursor homogeneity is deliberately examined in one study, solid-state and sol--gel preparation are compared in another, and a solid-state route is followed through lithium and powder-geometry controls in a third \cite{parascos2022compositional,kosir2022comparative,heywood2023tailoring}. The present records do not establish a common-powder, matched-composition, matched-thermal-budget comparison with a shared replication protocol. Route chemistry, mixing scale, dopant distribution, and later firing may therefore explain the same apparent pattern.

A discriminating study would use one reagent lot and a predeclared factorial design in which mixing or powder route changes while nominal composition, dopant, green-body preparation, and final thermal schedule are held fixed. Independent batches would be assessed with the same phase, density, and separated transport protocol. The working hypothesis that greater precursor homogeneity reduces outcome dispersion would be falsified if the prespecified dispersion measure is not lower than the matched solid-state control. This action tests an explanation; it does not assume the conclusion.

\subsection{Gap 2: starting lithium excess versus firing-stage lithium control}

The abstracts for a solid-state route and a Ga-containing processing study report lithium inventory among their controlled variables, while the cover-powder abstract reports a lithium-rich firing boundary together with density and transport readouts \cite{heywood2023tailoring,schwab2023towards,zhang2019densification}. Across those records, starting excess, dwell, specimen surface-to-volume ratio, and cover-powder supply are not isolated on one common composition.

A predeclared two-by-two design can separate nominal lithium excess (present/absent) from lithium-supplying cover material (present/absent), with common powder, crucible, geometry, and thermal history. The record should close a lithium balance as far as the selected analytical method permits and retain phase, density, and transport uncertainty. An interaction explanation would be weakened or falsified if its preregistered interaction term were indistinguishable from zero within the study's declared uncertainty. The gap is the absence of this discrimination in the checked collection, not an assertion of field-wide absence.

\subsection{Gap 3: additive-enabled densification versus boundary-phase mediation}

Abstracts for SiO2, Li5AlO4, and AlN interventions each report multiple structural or electrochemical outcomes and include interpretations involving carbonate removal, lithium supply, or boundary products \cite{zhang2020enhanced,zhou2024li5alo4,zhang2022high}. Because additive identity, dose, lithium activity, reaction products, final density, and boundary response change together, these abstracts cannot determine whether transport differences are mediated by density, boundary chemistry, or both.

A falsifiable next action is an additive-free control plus a dose series from one parent powder, with density matched where feasible and phase assemblage, boundary composition, and total/bulk/grain-boundary response measured before cell tests. A proposed boundary-phase mediation would be falsified if its predeclared boundary signature were absent or failed to covary with the separated grain-boundary response after density was accounted for. Null and contradictory outcomes would remain part of the report.

\subsection{Gap 4: fair comparison of pressureless, pressure-assisted, and ultrafast firing}

One abstract compares solid-state sintering, SPS, and SPS plus heat treatment, and another links a hot-pressing series to density and transport partition \cite{xue2020spark,david2015microstructure}. Ultrafast high-temperature processing is represented locally only by a title-level record \cite{ihrig2021li7la3zr2o12}. These provenance levels cannot support a cross-paper ranking, especially when starting powder, dopant, sample geometry, applied pressure, peak temperature, dwell, post-anneal, and conductivity temperature may differ.

A common characterized powder batch could be divided among pressureless, hot-press/SPS, and ultrafast protocols. Target density and specimen geometry would be declared in advance, energy/time exposure recorded, and one impedance protocol applied. The working equivalence hypothesis---that acceleration preserves transport at matched density---would be falsified if the accelerated route fell outside a prespecified equivalence margin relative to the pressureless control. This design compares routes on declared endpoints instead of method names.

\subsection{Gap 5: exposure history as a reproducibility variable}

Crucible and air-exposure outcomes are jointly reported in one stored abstract, lithium--proton exchange is presented as a processing-control variable in another, and precursor homogeneity is linked to reproducible material characteristics in a third \cite{xia2016ionic,rosen2021controlling,parascos2022compositional}. The current collection does not expose humidity, duration, storage state, surface carbonate or proton-exchange indicator, and pre-measurement treatment as one shared reporting block. An apparent route or batch effect could therefore contain an uncontrolled exposure component.

Replicate pellets from one batch could be assigned to predeclared humidity--time histories, with a surface/compositional indicator measured before and after exposure and a common transport protocol repeated. The exposure hypothesis would be falsified if neither the selected surface-state metric nor transport changed systematically with cumulative exposure under the preregistered analysis. Reporting the null result would be as informative as reporting degradation.

\subsection{Gap 6: an interlaboratory minimum record for density and conductivity}

Abstract-bearing studies in the collection jointly expose particle-size and uniformity variables, phase/density/transport fields, and route/dopant comparisons \cite{guo2023uniform,yan2020submicron,kosir2022comparative}. They do not themselves constitute an interlaboratory round robin. Relative-density basis, spatial sampling, conductivity partition, test temperature, electrodes, replicate count, and uncertainty are not uniformly visible in the locally checked evidence.

A blind round robin on shared pellets and shared powder would permit each laboratory's local protocol to be compared with a common minimum record. The preregistered analysis would partition dispersion associated with specimen, measurement, and reporting choices. The hypothesis that harmonized reporting and analysis explain a material fraction of between-laboratory dispersion would be falsified if the relevant variance component did not decrease. The immediate output would be a comparability result, not a claim that one laboratory or method is superior.

Across all six gaps, the common priority is controlled discrimination. The process map identifies where a variable enters; the minimum record preserves the conditions needed to interpret it; and the falsification clause states what evidence would change the working explanation. That combination avoids converting an evidence limitation into either a novelty claim or an experimentally actionable recipe.
